\subsection{Determinante und Lösungstyp}
\label{subsec:linear_systems_determinant_solution_criterion}

Die Cramersche Regel scheint zunächst für jedes $n\times n$-LGS eine explizite
Lösung anzugeben. Das wirkt überraschend: Wir kennen bereits quadratische LGS
ohne eindeutige Lösung. Wie passt beides zusammen?

\begin{KeyPointBox}{Wann kann die Cramersche Regel scheitern?}{linear_systems_cramers_rule_denominator}
  In jeder Cramerschen Lösungsformel
  \[
    x_j=\frac{\det(A_j)}{\det(A)}
  \]
  steht derselbe Nenner $\det(A)$. Ist $\det(A)\neq0$, sind alle Quotienten
  definiert und Cramer liefert die eindeutige Lösung. Die Formel kann daher
  nur dann scheitern, wenn
  \[
    \det(A)=0
  \]
  ist. Die folgende Verbindung mit dem Rang zeigt, dass dies nicht nur ein
  technisches Versagen der Formel ist: Genau dann ist eine eindeutige Lösung
  unmöglich.
\end{KeyPointBox}

Die Determinante prüft die Eindeutigkeit somit, ohne dass das LGS vollständig
gelöst werden muss.

\begin{TheoremBox}{Determinantenkriterium für $n\times n$-LGS}{linear_systems_determinant_solution_criterion}
  Für ein $n\times n$-LGS mit Koeffizientenmatrix $A$ gilt
  \[
    \det(A)\neq0
    \quad\Longleftrightarrow\quad
    \operatorname{rang}(A)=n
    \quad\Longleftrightarrow\quad
    \text{genau eine Lösung}.
  \]
  Entsprechend gilt
  \[
    \det(A)=0
    \quad\Longleftrightarrow\quad
    \operatorname{rang}(A)<n
    \quad\Longleftrightarrow\quad
    \text{keine oder unendlich viele Lösungen}.
  \]
\end{TheoremBox}

\begin{ExampleBox}{Gleiche Determinante, verschiedene Lösungstypen}{linear_systems_zero_determinant_solution_types}
  Die beiden LGS besitzen dieselbe Koeffizientenmatrix
  \[
    A=\begin{pmatrix}1&1\\1&1\end{pmatrix},
    \qquad \det(A)=1\cdot1-1\cdot1=0,
  \]
  aber verschiedene Absolutglieder:

  \smallskip
  \noindent
  \begin{minipage}[t]{0.47\linewidth}
    \vspace{0pt}\centering
    \textbf{LGS I}\par\smallskip
    \(
      \begin{aligned}
        x_1+x_2&=2,\\
        x_1+x_2&=1
      \end{aligned}
    \)
    \par\medskip
    \Cloze{4.2cm}{\text{keine Lösung}}
  \end{minipage}\hfill
  \begin{minipage}[t]{0.47\linewidth}
    \vspace{0pt}\centering
    \textbf{LGS II}\par\smallskip
    \(
      \begin{aligned}
        x_1+x_2&=2,\\
        x_1+x_2&=2
      \end{aligned}
    \)
    \par\medskip
    \Cloze{5.0cm}{\infty\text{-viele Lösungen}}
  \end{minipage}

  \medskip
  Die Determinante erkennt in beiden Fällen, dass die Lösung nicht eindeutig
  ist. Zwischen keiner und unendlich vielen Lösungen entscheidet erst der
  Rangvergleich der erweiterten Koeffizientenmatrix.
\end{ExampleBox}

\begin{KeyPointBox}{Spezialfall: homogenes $n\times n$-LGS}{linear_systems_homogeneous_determinant_criterion}
  Ein homogenes LGS besitzt immer die triviale Lösung. Deshalb bleiben nur
  zwei Fälle:
  \[
    \begin{aligned}
      \det(A)\neq0
      &\quad\Longleftrightarrow\quad
      \text{nur die triviale Lösung},\\[0.8em]
      \det(A)=0
      &\quad\Longleftrightarrow\quad
      \text{unendlich viele Lösungen, darunter nichttriviale}.
    \end{aligned}
  \]
  Dieses Kriterium verwenden wir später, um lineare Abhängigkeit und
  Unabhängigkeit von Vektoren zu prüfen.
\end{KeyPointBox}
